% ============================================================================
% ch5-conclusions.tex — Chapter IV: Conclusions
% ============================================================================
\chapterpage{Chapter\kern0.25em IV}{Conclusions}{images/IM4.jpeg}
\thispagestyle{chapterstart}

\subsection*{Key Messages}
\greensep

\begin{twocolbody}

\subsubsection*{Key message 1 --- The scale of exposure}

The installed base of 8,572 data centers carries US\$388 billion in climate-related value at risk---38\% of enterprise value. With prospective AI infrastructure, total exposure rises to \$1.0--3.7 trillion depending on growth scenario. This is the first sector-wide quantification of climate risk as enterprise value loss for digital infrastructure. The gap between 38\% enterprise value erosion and zero in standard DCF valuations quantifies the mispricing. Correcting it requires a metric. ClimVaR provides one.

\subsubsection*{Key message 2 --- Geography determines the risk, not scale}

A twofold dispersion in risk intensity (28\% in the US to 56\% in China) and distinct channel signatures by region mean that a single global adaptation strategy systematically underinvests in the dominant risk channel. Business interruption dominates in the US (44\%), carbon pricing in Europe (43\%), physical damage in Asia-Pacific (39\%), and China carries compound exposure across three channels. Two identical facilities, operated by the same company, face different climate-adjusted valuations depending on where they sit. Channel attribution is not an analytical refinement---it is the precondition for any adaptation strategy that does not systematically misallocate capital.

\columnbreak

\subsubsection*{Key message 3 --- AI infrastructure inverts the risk profile}

AI-era infrastructure does not replicate the installed base's risk profile---it amplifies it. On the installed base, carbon costs represent 28\% of ClimVaR and physical channels 72\%. On AI-tier facilities, physical channels account for 94\% while carbon drops to 6\%. Net-zero commitments address the carbon component but leave 82--88\% of AI-era risk untouched. Each additional GW of AI capacity amplifies physical exposure through three mechanisms: geographic concentration in climate-vulnerable locations, higher power density per rack, and supply chain dependence on a narrow set of component suppliers. Sustainability strategies centered on renewable procurement address less than 10\% of this exposure. The other 90\% requires a different strategy: cooling architecture, geographic diversification, and supply chain redundancy.

\subsubsection*{Key message 4 --- Adaptation pays for itself}

Proactive adaptation reduces aggregate ClimVaR by 30--39\% with positive net present value under every scenario tested. All ten adaptation actions fall below the break-even line. Carbon is the most responsive channel ($-$63\% via PPAs). But the residual 61\% defines the limit of facility-level action: extreme physical events, supply chain concentration (geographically concentrated semiconductor fabrication, power transformer lead times), and thermodynamic limits on cooling in the hottest climates cannot be engineered away. Every year of delayed adaptation converts greenfield opportunity into future retrofit cost, permanently.

\end{twocolbody}

% ---- STAKEHOLDER IMPLICATIONS ----
\vspace{4mm}

\subsection*{Implications for Stakeholders}
\greensep

\begin{twocolbody}

\subsubsection*{Operators and system planners}

ClimVaR provides a metric compatible with the DCF models that govern investment decisions. The $2\times$ dispersion between regions makes site selection a risk management decision, not merely a cost optimization. The Shapley decomposition identifies the highest-return adaptation lever at each site: identical capital deployed at Site~A/US (supply chain dominated), Site~B/France (flood trapped), and Site~C/China (triple threat) generates fundamentally different risk reductions---invisible without channel attribution. Every data center built today embeds assumptions about climate futures. The framework makes those assumptions explicit and their financial consequences quantifiable.

\subsubsection*{Investors and insurers}

The installed base's loss experience understates the prospective portfolio's risk profile. The $\times$2.6 physical risk amplification on AI infrastructure and the $\times$11 supply chain multiplier explain why facility-level loss ratios underestimate systemic exposure---at every site examined, 98--99\% of business interruption originates upstream, not on site. The Shapley + Leontief decomposition offers a structure for differentiating premiums by channel composition rather than by geography alone. A portfolio with disclosed ClimVaR of 38\% should command a different price from one at 15\%---yet no disclosure framework currently enables this distinction. Building that framework is now a competitive advantage for the institutions that move first.

\subsubsection*{Policymakers and regulators}

Data centers have shifted from commercial real estate to strategic infrastructure, but their climate exposure remains weakly integrated into permitting, prudential oversight, and critical-infrastructure policy. The 38\% intensity on over \$1~trillion of installed assets exceeds the scope of most national infrastructure protection programs. The concentration of AI investment in climate-vulnerable locations---the US South, the US East Coast, southern China, Southeast Asia---represents a market failure requiring climate-adjusted siting rules, infrastructure stress tests, and public risk-sharing mechanisms calibrated to the actual financial exposure. The analytical basis for those mechanisms now exists. Integrating it into regulatory practice is the next step.

\subsubsection*{AI developers}

The sustainability agenda for AI infrastructure has been built around the wrong constraint. Physical resilience, not energy sourcing, is what determines whether a \$10 billion training cluster survives a climate event. The structural inversion from carbon-dominated to physically-dominated risk means that sustainability strategies focused on renewable procurement address less than 10\% of AI-era climate exposure. Cooling architecture, geographic diversification, supply chain redundancy, and operational continuity planning are the levers that matter most for protecting multi-billion dollar training investments from climate disruption.

\subsubsection*{Chief sustainability officers}

Climate risk reporting that stops at carbon emissions misses three quarters of the exposure. CSOs overseeing data center portfolios face a disclosure gap: current frameworks (CDP, TCFD, CSRD) emphasize transition risk metrics---Scope~1/2/3 emissions, renewable energy share, carbon intensity---while physical risk remains unquantified. The ClimVaR framework provides the missing metric. Integrating it into sustainability reporting would align disclosed risk with actual financial exposure---and position early adopters ahead of the disclosure curve, before regulators make the choice for them.

\end{twocolbody}

% ---- LIMITATIONS ----
\subsection*{Limitations}
\greensep

\begin{twocolbody}

\noindent The \$388~billion figure represents a lower bound on total economic exposure. Several factors not captured in the model would increase it.

\subsubsection*{Climate projection uncertainty}

Physical climate projections carry inherent uncertainty at facility-specific spatial scales. The scenario structure (RCP~2.6/8.5 $\times$ NZT/BAU) spans a wide portion of the assessed range, but tail risks---tipping point cascades, correlated infrastructure failures, compound events exceeding historical analogs---may exceed modeled bounds. The January 2026 US East Coast winter storm illustrates the type of correlated event the framework captures directionally but cannot predict with precision.

\subsubsection*{Archetype-based calibration}

All facilities within a region and tier share financial parameters (PUE, PPA factor, WACC, cost ratios), reflecting industry benchmarks from JLL, Synergy Research, CBRE, and IEA. This does not capture facility-specific characteristics: construction materials, elevation, cooling system vintage, backup generator capacity. The framework is most reliable for portfolio-level and regional comparison; individual facility estimates carry uncertainty proportional to their divergence from the archetype.

\subsubsection*{Prospective infrastructure}

The prospective facilities (615--2,939 depending on growth scenario) are based on announced, planned, and speculative projects. Actual buildout will differ in timing, scale, and location. ClimVaR estimates for prospective infrastructure should be read as indicative of the risk profile under each scenario, not as point predictions.

\subsubsection*{Adaptation modeling}

The Selective and Proactive plans apply uniform reduction rates by channel. In practice, adaptation effectiveness varies by facility age, operator capability, and local regulatory context. The $\times$2.6 amplification finding combines genuine physical amplification with a valuation effect (AI facilities carry higher asset value per MW); separating these two components matters for adaptation strategy, since actions that reduce physical exposure differ from those that address valuation premiums.

\subsubsection*{Scope of risk}

The framework quantifies private financial risk through five channels. It does not capture systemic externalities: grid strain from concentrated AI load, water resource competition, community displacement, or cascading economic effects of simultaneous service interruptions across interconnected facilities.

Nor does it capture several risk categories that are likely to become material over the asset lifetime. \textbf{Nature-related risk}: water stress, biodiversity constraints on siting permits, and ecosystem degradation around cooling water sources. \textbf{Reputation and litigation risk}: public opposition to data center siting, environmental litigation, and greenwashing claims against operators whose sustainability reporting omits physical exposure. \textbf{Technology risk}: obsolescence of cooling architectures, stranded assets from chip architecture shifts, and lock-in effects from early design choices. \textbf{Investor sentiment risk}: portfolio-level repricing as climate disclosure frameworks mature and physical risk metrics become standardized. Firm-level ClimVaR outputs will likely remain proprietary---granular site valuations carry material disclosure risk and competitive sensitivity. Yet capital markets will progressively internalize physical climate exposure through rising risk premia, rating adjustments, and stranded-asset repricing, regardless of voluntary disclosure. The question is not whether repricing occurs, but whether it proceeds in an orderly or disruptive fashion. \textbf{Electricity price risk}: the feedback loop between data center demand growth, grid congestion, and wholesale electricity prices---particularly in markets where AI load growth outpaces generation capacity. Each of these channels would increase, not decrease, the total exposure estimate.

\end{twocolbody}

\subsection*{Future Research}
\greensep

\begin{twocolbody}

\noindent Several extensions would strengthen and generalize the analysis.

\subsubsection*{Empirical calibration}

The January 2026 US East Coast winter storm provides a natural experiment for calibrating the $\beta$ channel against observed losses. Systematic collection of post-event operational and financial data---downtime duration, revenue impact, insurance claims, recovery costs---would enable model validation against real-world outcomes. Collaboration with operators and insurers would be required.

\subsubsection*{Water-energy-climate nexus}

Liquid cooling reduces energy consumption but increases water consumption. In water-stressed regions---the US Southwest, parts of India, northern China---this substitution may exchange one climate vulnerability for another. Integrating water availability projections with cooling technology trajectories would close a gap in the current analysis and is particularly relevant as AI facilities adopt liquid cooling at scale.

\subsubsection*{Real options for adaptation timing}

The current framework treats adaptation as a binary decision (invest or not). A real-option extension would value the flexibility to delay while preserving the ability to act---relevant for legacy facilities where retrofit timing is uncertain and technology is evolving rapidly. The 15--20\% versus 40--60\% CAPEX differential between greenfield and retrofit suggests significant option value in early commitment.

\subsubsection*{Expected Shortfall and tail risk metrics}

The current analysis reports ClimVaR at Q25, MEAN, and Q95 percentiles. Extending this to Expected Shortfall (CVaR)---the conditional expectation of losses beyond a given quantile---would provide a risk metric better suited to insurance pricing and prudential regulation, where tail losses dominate decision-making. The mathematical foundations are developed in Appendix~A.

\subsubsection*{Cross-sector generalization}

The combination of entity-level ClimVaR, Shapley channel attribution, and Leontief supply chain propagation provides a transferable methodology. Manufacturing, energy, and transportation infrastructure face analogous challenges of long-lived capital assets under climate uncertainty. Extension to multi-sector infrastructure portfolios---particularly at points of interdependency (data centers $\times$ power grids $\times$ telecommunications)---would capture systemic risks that sector-specific analysis underestimates.

\subsubsection*{System Dynamics Modeling}

The current model compares outcomes across fixed scenario endpoints. A natural extension would make the pathways themselves time-varying: carbon prices responding to policy signals, operator strategies adapting to observed losses, and MRIOT coefficients reflecting supply chain restructuring over time. Dynamic multi-regional input-output modeling along these lines is being developed in parallel work \citep{beaufils2025}.

\end{twocolbody}

\vspace{6mm}

\begin{highlightbox}
The prospective infrastructure whose design has not been finalized---between 63 and 296~GW depending on growth scenario---represents the remaining degrees of freedom. The analytical basis for informing those decisions now exists. Its integration into investment practice is the next step.
\end{highlightbox}
